An essential requirement for relating logical corpora is standardizing the identifiers so that each identifier in the corpus can be uniquely referenced.
It is desirable to use a uniform naming schema so that the syntax and semantics of identifiers can be understood and implemented as generically as possible.
Therefore, we use MMT URIs \cite{RK:mmt:10}, which have been specifically designed for that purpose.

\subsubsection{Namespace Organization}

MMT URIs standardize the syntax of the identifiers, but they still allow a lot of freedom how to assign URIs to the concepts in a specific corpus.
This assignment is straightforward in principle --- after all we only have to make sure that every concept has a unique URI.
However, as we will see below, the structure of a corpus can pose some subtle issues that must be addressed carefully. 
Therefore, we quickly discuss commonly used corpus structures and how these can be used to form URIs systematically.

The common structuring feature of corpora is usually a directory tree.
The leaves of this tree are files and contain modules.
Moreover, each corpus usually has a certain root namespace.
However, systems differ in how they subdivide a corpus into namespaces.

We distinguish the following cases:
\begin{compactitem}
 \item \textbf{flat} structure: All files share the same namespace regardless of their physical location in the directory tree.
 This naming schema is most well-known from SML.
 In this case, we can use the root namespace as the fixed namespace for all concepts in the corpus.
 \item \textbf{directory-based} structure: The namespace of a module is formed by concatenating the root namespace with the path to the directory containing it.
 There are two subcases regarding the treatment of the file name:
  \begin{compactitem}
    \item \textbf{files-as-modules}: Each file contains exactly one module, whose name is that of the file without the file name extension.
    The name of the module may be repeated \emph{explicit}ly in the file or may be left \emph{implicit}.
    Files as explicitly named modules is most well-known as the convention of Java.
    \item \textbf{irrelevant file names}: The file name is irrelevant, i.e., the grouping of modules into files within the same directory is arbitrary.
    In particular, a file can contain multiple modules.
  \end{compactitem}
 \item \textbf{file-based} structure: The namespace of a module is formed by concatenating the root namespace of the corpus with the path to the file containing it.
\end{compactitem}

\subsubsection{URIs for Selected Proof Assistants}

Using the principles defined above, we describe the MMT URI formation principles for some important proof assistants.
In all cases, we also assign MMT URIs for the underlying foundations in order to refer to built-in concepts.
\medskip

\textbf{PVS} \cite{pvs} uses directory-based namespaces with irrelevant file names.
We propose the following root URIs for some important PVS-related corpora:
\begin{compactitem}\small
 \item the PVS foundation: \url{http://pvs.csl.sri.com/foundation}
 \item the standard library that is shipped with PVS: \url{http://pvs.csl.sri.com/Prelude}
 \item the NASA corpus: \url{http://shemesh.larc.nasa.gov/fm/ftp/larc/PVS-library}
\end{compactitem}
Within a PVS corpus, the top-level modules are theories and (co)datatype declarations.
The only possible nesting between them is that theories may contain (co)datatype
declarations.  Consequently, the module names have at most two segments.  If a module is a
(co)datatype, its symbols are the constructors, testers, etc.  These are not hierarchical.
If a module is a theory, its symbols are all declarations declared in it.  These are not
hierarchical.  However, if a symbol name $N$ is declared multiple times in the same module
(due to overloading), we use two-level names of the form $N/i$ where $i$ numbers all
declarations of $N$ in that module (starting from $1$).  \medskip

\textbf{Coq} \cite{coq} uses directory-based namespaces with files as implicitly named modules.
We propose the following root URIs:
\begin{compactitem}
 \item the Coq foundation: \url{https://coq.inria.fr/foundation}
 \item the standard library shipped with Coq: \url{https://coq.inria.fr/theories}
 \item the Coq contributions: \url{https://coq.inria.fr/contribs}
 \item the Mathematical Components corpus (including SSReflect):\\ \url{http://ssr.msr-inria.inria.fr/math-comp}
\end{compactitem}
Coq modules can be nested.
Besides the module name given implicitly by the file, Coq files can contain modules and module types.
Symbols are all declarations inside a module.
Their names can be hierarchic due to generative functor instantiation.
\medskip

\textbf{Matita} uses the same URIs as Coq except for not allowing nested modules.
We suggest the following root URIs:
\begin{compactitem}
\item the Matita foundation: \url{http://matita.cs.unibo.it/foundation}
\item the Matita standard library: \url{http://matita.cs.unibo.it/library}
\end{compactitem}
\medskip

\textbf{Mizar} uses a flat namespace.
We propose the following root URIs:
\begin{compactitem}
 \item the Mizar foundation: \url{http://mizar.org/foundation}
 \item the Mizar Mathematical Library (MML): \url{http://mizar.org/library}
\end{compactitem}

The Mizar modules are the articles.
The name of a module is the name of the article without the file name extension.
There is no nesting of modules.
%
The Mizar symbols are all the declarations inside an article.
Their names are obtained through a heavily idiosyncratic naming schema that includes generating unique names by numbering the declarations of the same kind in each article.
For example, the MMT URI of conjunction is
\url{http://mizar.org/corpus?XBOOLEAN0?K4}
\medskip

\textbf{HOL Light} does not have an obvious MMT URI formation principle because it does not maintain all its identifiers itself --- instead it relies on the OCaml toplevel to store the assigned values.
There are three kinds of HOL Light symbols: types, constants, and theorems; only the former two are visible to the HOL Light kernel.
For each type or constant, we use the name visible to the HOL Light kernel.
For each theorem, we use the OCaml binding name.
OCaml toplevel symbols may be grouped into OCaml modules. This feature is seldom used, as it only affects theorem names: the kernel is also not aware of the OCaml module in which the definitions of constants or types are introduced.

This has the effect that different HOL Light files can be incompatible with each other.
There are two reasons for this incompatibility: First, toplevel symbols may be overwritten by others, which means that the original ones are
no longer accessible and a formalization might fail.
This happens for example in case of the OCaml basic output function \texttt{open\_in}
which gets overwritten by a theorem with the same name.
Second, the loading of a file may change the state of HOL Light kernel or packages to one that is no longer compatible
with another file~\cite{wiedijk-types09}. For example the theories \texttt{complexnumbers.ml}
and \texttt{complexes.ml} both introduce the type \texttt{complex} but with different
definitions.

Thus, HOL Light's names do not uniquely identify symbols.  Therefore, we use
directory-based namespaces with files-as-modules.  For constants and types introduced by a
module we add the prefixes \texttt{const/} and \texttt{type/} respectively.  If a file
contains OCaml modules, we use their names to form multi-segment module names.
Accordingly, if symbols result from OCaml structures, we form multi-segment symbol names.
This has the effect that HOL Light URIs are formed in exactly the same way as for Coq.

We propose the following root URIs:
\begin{compactitem}
 \item the HOL Light foundation and the library shipped with it: \\ \url{http://github.com/jrh13/hol-light}
 \item the Formal Proof of Kepler formalization:\\
   \url{http://github.com/flyspeck/flyspeck}
\end{compactitem}

The theorem \texttt{well\_defined\_unordered\_pair} from Flyspeck~\cite{hales-pisigma15} is for example assigned the complete
URI: \texttt{http://github.com/flyspeck/flyspeck ? text\_formalization/packing/marchal3 ? Marchal\\\_cells\_3.well\_defined\_unordered\_pair}.
\medskip

\textbf{HOL4} internal module names correspond to names of files, however the names
of types and constants are not associated with the modules. Furthermore, the names of
constants and types are separate. To avoid ambiguity we use the same module names as
for HOL Light.% We propose the following root URIs:

\begin{compactitem}
 \item the HOL4 foundation and the library shipped with it:\\
 \url{https://hol-theorem-prover.org}
\end{compactitem}
\medskip

\textbf{Isabelle} is a logical framework: Its distribution includes a number of object logics, and each Isabelle theory uses an object logic (or declares a new one).
Isabelle uses files as explicitly named modules.
However, it disregards the directory structure: The system makes sure that two modules with same name cannot be loaded in the same session even if they are stored in different directories.
But as different object logics and different developments often declare incompatible notions, we still use directory-based namespaces to make sure all theories have unique namespaces.

Isabelle allows several module mechanisms including locales and type classes
\cite{haftmannw-types06}.  Therefore, we form nested module names by concatenating theory
and locale/type class names. % We propose the following root URIs:
\begin{compactitem}
 \item Isabelle foundation and distributed libraries: \url{http://isabelle.in.tum.de/}
 \item the Archive of Formal Proofs: \url{http://afp.sf.net/}
 \item the TLA+ logic: \url{http://tla.msr-inria.inria.fr/}
 seL4: https://github.com/seL4
\end{compactitem}

For example the type of streams is represented by the URI:\\
\url{http://isabelle.in.tum.de/ ? HOL/corpus/Stream ? stream}


For informal collections of mathematical knowledge like Wikipedia, we can usually adopt
similar measures. We interpret Wikpedia as having a flat namespace
\url{http://en.wikipedia.org/wiki}, obtain modules via the files-as-models regime, and use
the anchors of editable fragments as symbol names. So for instance the statement of
uniqueness of identity element and inverses in the Wikipedia article on groups would have
the MMT URI
\url{https://en.wikipedia.org/wiki?Group_(mathematics)?Uniqueness_of_identity_element_and_inverses}.


%%% Local Variables:
%%% mode: latex
%%% TeX-master: "report"
%%% End:

%  LocalWords:  standardizing mathtt mathtt mathtt compactitem textbf emph emph medskip
%  LocalWords:  pvs Mizar ednote formalization texttt wiedijk-types09 complexnumbers.ml
%  LocalWords:  hales-pisigma15 haftmannw-types06 openaxiom
